% capitulo-3.tex
% Capítulo 3 — Metodología
%
% This chapter describes the planned experimental methodology. Citations
% use the citation keys of the project Zotero library (exported to
% references.bib). Items that depend on decisions not yet finalized
% (hardware details, software versions, model selection, resolutions,
% durations, repetitions) are marked with visible \todo{...} placeholders
% to be completed once determined; % TODO comments mark citation gaps. No
% experimental results are reported in this chapter.

\providecommand{\todo}[1]{\textbf{[TODO: #1]}}

\chapter{Metodología}
\label{chap:methodology}

\section{Plataforma experimental}
\label{sec:platform}

La plataforma experimental principal es una NVIDIA Jetson Orin Nano,
elegida como módulo embebido representativo para la visión por computador
en el borde \cite{scalcon2024, alqahtani2026}. Las siguientes
características del dispositivo se registran y se mantienen constantes a lo
largo de los experimentos:

\begin{itemize}
  \item Módulo/modelo de Jetson: \todo{SPECIFY: exact module and developer
        kit version}
  \item Capacidad de RAM: \todo{SPECIFY}
  \item Almacenamiento: \todo{SPECIFY}
  \item Modo de energía: \todo{SPECIFY: fixed power mode for all
        experiments. Published Orin-family evaluations distinguish power
        modes when reporting frame rates \cite{scalcon2024}, so this
        setting is part of the documented experimental condition.}
  \item Configuración de enfriamiento: \todo{SPECIFY. Sustained inference
        load raises device temperature \cite{zequn2026}, so cooling is
        fixed and documented.}
  \item Cámara / dispositivo de entrada conectado: \todo{SPECIFY}
  \item Temperatura de operación, donde corresponda: \todo{SPECIFY: how
        thermal conditions are monitored}
\end{itemize}

Todos los experimentos se ejecutan en este único dispositivo, con los nodos
de las configuraciones de ROS~2 ejecutándose como procesos separados del
sistema operativo en el mismo anfitrión. Ejecutar todos los procesos en un
dispositivo aísla el efecto de la propia arquitectura de ROS~2 de los
efectos de red, que están fuera del alcance de esta tesis.

\section{Entorno de software}
\label{sec:software-environment}

El entorno de software completo se documenta y se mantiene constante:

\begin{itemize}
  \item Sistema operativo: \todo{SPECIFY: Linux distribution and version}
  \item Versión de JetPack: \todo{SPECIFY}
  \item Distribución de ROS~2: \todo{SPECIFY: one distribution, fixed for
        all experiments}
  \item Implementación de DDS: \todo{SPECIFY: one implementation, fixed
        for all experiments, selected through the rmw layer}
  \item Configuración de QoS: \todo{SPECIFY: one fixed configuration,
        stated explicitly, including reliability, durability, and history
        policies. QoS settings measurably affect latency
        \cite{fernandez2020, kronauer2021} and must therefore be part of
        the documented experimental condition.}
  \item Versión de Python: \todo{SPECIFY}
  \item Versiones de paquetes (entorno de ejecución de inferencia,
        bibliotecas de imágenes, herramientas de monitoreo): \todo{SPECIFY}
\end{itemize}

Para los experimentos principales se utiliza un despliegue sin
contenedores. Las decisiones anteriores se tratan como condiciones
experimentales controladas y no como variables bajo investigación; cada una
se documenta de modo que los experimentos puedan reproducirse
(Sección~\ref{sec:reproducibility}).

\section{Cadena de procesamiento de visión por computador}
\label{sec:pipeline}

La carga de trabajo es una cadena de procesamiento de detección de objetos
que consiste en las siguientes etapas lógicas:

\begin{enumerate}
  \item \textbf{Adquisición de imágenes:} obtener cada fotograma de
        entrada.
  \item \textbf{Preprocesamiento:} convertir el fotograma al formato
        esperado por el modelo (p. ej., redimensionamiento, conversión de
        color, normalización y disposición del tensor). Esta etapa es
        individualmente significativa: para la detección de objetos, el
        preprocesamiento contribuye una porción sustancial del tiempo de
        ejecución total, especialmente a resoluciones más altas
        \cite{szegedi2026}.
  \item \textbf{Inferencia:} ejecutar el modelo de red neuronal.
  \item \textbf{Postprocesamiento:} convertir la salida sin procesar del
        modelo en resultados a nivel de aplicación (p. ej., cuadros
        delimitadores, filtrado por confianza, supresión de no máximos). Al
        igual que el preprocesamiento, esta etapa no es despreciable para
        la detección de objetos \cite{szegedi2026}.
  \item \textbf{Salida:} producir el resultado final de la cadena de
        procesamiento.
\end{enumerate}

El modelo de red neuronal se selecciona de acuerdo con los criterios de la
Sección~\ref{sec:workload}, sobre la base de evidencia publicada de que la
detección de objetos en tiempo real es factible en la Orin Nano con las
familias estándar de detectores de una etapa \cite{scalcon2024,
alqahtani2026, szegedi2026,zequn2026}:
\todo{SPECIFY: final model, weights version, precision, and inference
runtime once selected.} El mismo modelo y la misma configuración de
inferencia se utilizan en cada configuración experimental.

\section{Implementación monolítica}
\label{sec:monolithic}

La implementación de línea base ejecuta la cadena de procesamiento completa
dentro de un único proceso de Python:

\begin{quote}
\texttt{Entrada $\rightarrow$ Preprocesamiento $\rightarrow$ Inferencia
$\rightarrow$ Postprocesamiento $\rightarrow$ Salida}
\end{quote}

Todas las etapas se ejecutan secuencialmente dentro de un proceso. Los datos
pasan entre las etapas como estructuras de datos en memoria, sin
serialización ni transporte entre procesos. Esta implementación sirve como
línea base de rendimiento contra la cual se comparan las implementaciones de
ROS~2. Corresponde al extremo arquitectónico en el que la aplicación
aprovecha plenamente el hecho de que todos los datos de la cadena pueden
permanecer en un único espacio de direcciones.

\section{Implementación en ROS 2}
\label{sec:ros2-implementation}

Las mismas etapas lógicas se implementan como nodos de ROS~2 utilizando
rclpy, cada uno ejecutándose como un proceso separado del sistema operativo
en el mismo anfitrión, conectados en cadena mediante tópicos que
transportan mensajes de imagen:

\begin{quote}
\texttt{Nodo A $\rightarrow$ Nodo B $\rightarrow$ Nodo C $\rightarrow$
Nodo D}
\end{quote}

Cada nodo publica y se suscribe con la configuración de QoS fija, y la
configuración del ejecutor es igualmente fija e idéntica entre los nodos,
dado que el ejecutor es inseparable del rendimiento del nodo
\cite{macenski2023, casini2026}. El mapeo de las etapas de procesamiento a
nodos preserva la equivalencia funcional con la implementación monolítica:
dada la misma secuencia de entrada y el mismo modelo, ambas implementaciones
producen salidas equivalentes.
\todo{SPECIFY: the exact stage-to-node mapping for each configuration once
finalized, the executor type used, and how functional equivalence is
verified (e.g., comparing outputs on a fixed input sequence).}

Esta implementación corresponde a la manera ordinaria y recomendada por
omisión de estructurar una aplicación de ROS~2 como nodos separados
\cite{macenski2022}, antes de que se aplique cualquiera de las mitigaciones
descritas en la literatura relacionada---composición, comunicación
intraproceso o transporte por memoria compartida \cite{macenski2023,
kamiya2025}---. Estudiar esta línea base de forma aislada es deliberado:
mide la arquitectura que una adopción directa de ROS~2 produce.

\section{Configuraciones experimentales}
\label{sec:configurations}

Se utiliza un número reducido de configuraciones, elegidas de modo que
representen diferencias arquitectónicas significativas y no divisiones
arbitrarias de la cadena de procesamiento:

\begin{itemize}
  \item \textbf{Configuración A --- Monolítica:} un proceso de Python.
  \item \textbf{Configuración B --- Particionamiento bajo:} dos nodos de
        ROS~2.
  \item \textbf{Configuración C --- Particionamiento mayor:} cuatro nodos
        de ROS~2.
\end{itemize}

\todo{SPECIFY: the exact node boundaries for Configurations B and C, and
the justification that they represent meaningful architectural
differences.}

La matriz experimental principal cruza las tres configuraciones
arquitectónicas con tres resoluciones de entrada. La resolución se incluye
como dimensión secundaria porque ambos términos del compromiso bajo estudio
escalan con ella: el costo de la inferencia crece con la resolución de
entrada \cite{szegedi2026}, y también el costo de comunicación de los
propios mensajes de imagen, dado que los costos de serialización y copia
por mensaje crecen con el tamaño de la carga útil \cite{kronauer2021,
macenski2022, luo2025}. Las resoluciones se seleccionan según las
capacidades de la fuente de entrada y del modelo elegidos:

\begin{itemize}
  \item Resolución 1: \todo{SPECIFY (candidate: 320$\times$240)}
  \item Resolución 2: \todo{SPECIFY (candidate: 640$\times$480)}
  \item Resolución 3: \todo{SPECIFY (candidate: 1280$\times$720)}
\end{itemize}

Esto da como resultado $3 \times 3 = 9$ condiciones experimentales:

% NOTE: requires \usepackage{amssymb} in the preamble for \checkmark.
\begin{table}[htbp]
  \centering
  \caption{Matriz experimental principal.}
  \label{tab:experimental-matrix}
  \begin{tabular}{clcccc}
    \hline
    \textbf{Config.} & \textbf{Arquitectura} & \textbf{Particionamiento} &
    \textbf{Res.\ 1} & \textbf{Res.\ 2} & \textbf{Res.\ 3} \\
    \hline
    A & Proceso único & 1 proceso & \checkmark & \checkmark & \checkmark \\
    B & ROS 2          & 2 nodos   & \checkmark & \checkmark & \checkmark \\
    C & ROS 2          & 4 nodos   & \checkmark & \checkmark & \checkmark \\
    \hline
  \end{tabular}
\end{table}

Cada condición se repite múltiples veces, o se ejecuta por una duración fija
suficientemente larga, para obtener mediciones estables.
\todo{SPECIFY: number of repetitions and run duration per condition.}

\section{Procedimiento de evaluación comparativa de rendimiento}
\label{sec:procedure}

Cada condición experimental sigue el mismo procedimiento:

\begin{enumerate}
  \item \textbf{Inicializar el sistema.} Iniciar la aplicación y todas las
        herramientas de monitoreo.
  \item \textbf{Calentamiento.} Descartar una cantidad inicial de
        fotogramas o un período fijo de calentamiento para evitar medir el
        comportamiento de inicio: los costos de la primera solicitud, como
        la carga del modelo y la asignación de memoria, no reflejan la
        operación en estado estacionario.
        \todo{SPECIFY: warm-up duration or frame count.}
  \item \textbf{Ejecutar la carga de trabajo.} Ejecutar la cadena de
        procesamiento durante una duración fija o una cantidad fija de
        fotogramas.
        \todo{SPECIFY.}
  \item \textbf{Registrar las mediciones.} Recopilar la latencia de extremo
        a extremo, FPS, la utilización de CPU, RAM y GPU, y el consumo de
        energía si se mide.
  \item \textbf{Repetir.} Repetir el experimento lo suficiente para
        caracterizar la variabilidad.
  \item \textbf{Reiniciar.} Devolver el sistema al mismo estado de
        referencia antes de la siguiente condición experimental.
        \todo{SPECIFY: the reset procedure (e.g., stopping all nodes,
        cooling interval).}
\end{enumerate}

A lo largo de las mediciones, el sistema se mantiene en un estado
controlado: se aplica el modo de energía fijo y se minimiza la carga en
segundo plano. Esto sigue la práctica de la literatura de medición de
ROS~2, donde el comportamiento adaptativo del sistema operativo---el
escalado de frecuencia de la CPU y los estados de inactividad--- se
neutraliza explícitamente porque agrega ruido a las mediciones de latencia
\cite{kronauer2021}.

Los experimentos piloto preceden a los experimentos finales con el fin de
identificar problemas de medición, cuellos de botella inesperados,
supuestos incorrectos y errores de configuración. El propósito de las
ejecuciones piloto es fijar la metodología, no generar resultados finales.

\section{Medición de las métricas de rendimiento}
\label{sec:metric-measurement}

\subsection{Latencia de extremo a extremo}

La latencia se mide por fotograma como el tiempo transcurrido entre la
entrada del fotograma a la primera etapa de la cadena y la disponibilidad de
su salida completamente postprocesada. Las marcas de tiempo se toman en
puntos equivalentes en todas las configuraciones---equivalentes a las
fronteras de la cadena como un todo, no de las etapas internas
individuales--- de modo que las mediciones sean comparables entre
arquitecturas. Esto refleja la perspectiva de extremo a extremo de las
evaluaciones de latencia de ROS~2 a nivel de aplicación \cite{betz2023a} y
la metodología de marcas de tiempo en fronteras utilizada para el análisis
de comunicación resuelto por capa \cite{kamiya2025}.
\todo{SPECIFY: timestamping mechanism (e.g., monotonic clocks) and, for
the ROS 2 configurations, how the timestamps of the first and final nodes
are related on the same host.} Estadísticas informadas: media, mediana,
desviación estándar, mínimo, máximo, P95 y P99 cuando haya muestras
suficientes, siguiendo la práctica de informe de percentiles de la literatura
de evaluación comparativa de rendimiento de ROS~2 \cite{kronauer2021,
macenski2022}.

\subsection{Tasa de procesamiento}

La tasa de procesamiento se mide en fotogramas por segundo sobre el
intervalo de medición en estado estacionario, tras eliminar el calentamiento:

\begin{equation}
  \mathrm{FPS} = \frac{\text{fotogramas procesados}}{\text{tiempo
  transcurrido}}
  \label{eq:fps}
\end{equation}

\subsection{Utilización de CPU}

La utilización promedio y máxima de la CPU se mide durante el intervalo de
medición utilizando herramientas de monitoreo del sistema.
\todo{SPECIFY: tool and sampling interval (e.g., \texttt{tegrastats} or
equivalent).}

\subsection{Utilización de memoria}

La utilización de la RAM del sistema se mide durante el intervalo de
medición; el uso de memoria de GPU se registra si se dispone de una medición
confiable.
\todo{SPECIFY: tool.}

\subsection{Utilización de GPU}

La utilización de GPU se registra para determinar si el sobrecosto
arquitectónico cambia la utilización del recurso de cómputo principal.
\todo{SPECIFY: tool.}

\subsection{Consumo de energía (opcional)}

Si se dispone de una medición confiable: potencia promedio, potencia
máxima, energía por fotograma procesado y rendimiento por vatio, siguiendo
la práctica de métricas de potencia de los estudios de caracterización de la
Orin Nano \cite{zequn2026}.
\todo{SPECIFY: measurement method, or state that power measurement was
excluded. Power measurement is an optional extension and must not
threaten completion of the thesis.}

\section{Análisis de datos}
\label{sec:analysis}

Para cada condición experimental, el análisis procede de la siguiente
manera:

\begin{enumerate}
  \item Eliminar las mediciones de inicio y de calentamiento.
  \item Calcular estadísticas descriptivas de la latencia, la tasa de
        procesamiento y la utilización de recursos.
  \item Comparar las distribuciones de latencia entre configuraciones.
  \item Comparar la tasa de procesamiento y la utilización de recursos
        entre configuraciones.
  \item Calcular las diferencias porcentuales entre configuraciones,
        incluido el sobrecosto relativo de cada configuración de ROS~2 con
        respecto a la línea base monolítica:
        \begin{equation}
          \mathrm{sobrecosto}\ (\%) = \frac{L_{\mathrm{ROS}} -
          L_{\mathrm{mono}}}{L_{\mathrm{mono}}} \times 100
          \label{eq:overhead}
        \end{equation}
        donde $L$ denota la métrica bajo comparación (p. ej., la latencia
        media de extremo a extremo).
  \item Investigar si el aumento del particionamiento de nodos produce una
        tendencia consistente a través de las resoluciones.
\end{enumerate}

Donde corresponda, las pruebas estadísticas se seleccionan tras examinar los
datos recopilados y su distribución.
\todo{SPECIFY: tests to be applied once the data distribution is known.}

El análisis evalúa las siguientes hipótesis, enunciadas aquí como
predicciones y no como expectativas que deban confirmarse; los resultados
que contradigan una o más hipótesis siguen siendo válidos y se analizarán
como tales:

% NOTE: The thesis plan lists the hypotheses in the proposal but does not
% assign them a dedicated section in the chapter structure; they are
% placed here because the analysis evaluates them. Move if a better
% location is chosen.
\begin{description}
  \item[H1 (Latencia):] Las implementaciones de ROS~2 exhibirán una mayor
        latencia de extremo a extremo que la implementación monolítica
        debido al sobrecosto de comunicación y procesamiento introducido
        entre los nodos.
  \item[H2 (Tasa de procesamiento):] La implementación monolítica
        alcanzará una tasa de procesamiento igual o mayor que las
        implementaciones de ROS~2 bajo condiciones computacionales
        equivalentes.
  \item[H3 (Utilización de recursos):] Las implementaciones de ROS~2
        requerirán recursos adicionales de CPU y memoria en comparación
        con la implementación monolítica.
  \item[H4 (Particionamiento):] El aumento del número de fronteras de
        comunicación de ROS~2 incrementará el sobrecosto de comunicación
        y, por lo tanto, la latencia de extremo a extremo.
\end{description}

Estas predicciones se fundamentan en la literatura relacionada: el
sobrecosto de comunicación entre procesos se ha medido repetidamente
\cite{kronauer2021, macenski2022, kamiya2025}, y las mitigaciones que
reducen la cantidad de fronteras de proceso recuperan recursos sustanciales
\cite{macenski2023}. Las hipótesis requieren, sin embargo, verificación
experimental para esta carga de trabajo, plataforma y biblioteca cliente,
dado que la magnitud del sobrecosto en relación con el cómputo dominante de
la inferencia es exactamente lo que los estudios previos no aíslan.

\section{Reproducibilidad}
\label{sec:reproducibility}

La siguiente información se documenta de modo que los experimentos puedan
reproducirse, en línea con el llamado de la comunidad a una evaluación
comparativa de rendimiento robótica estandarizada y reproducible
\cite{mayoral-vilches2024}:

\begin{itemize}
  \item Modelo y configuración de hardware
  \item Versión de JetPack
  \item Distribución de Linux
  \item Distribución de ROS~2
  \item Implementación de DDS
  \item Versión de Python
  \item Versiones de paquetes
  \item Versión del modelo y pesos
  \item Configuración y precisión del modelo
  \item Resolución de entrada por experimento
  \item Configuración de QoS de ROS~2
  \item Modo de energía
  \item Configuración de CPU/GPU
  \item Condiciones ambientales y del entorno, donde corresponda
  \item Duración del benchmark
  \item Cantidad de repeticiones
  \item Procedimiento de calentamiento
\end{itemize}

Registrar la configuración completa de software y hardware no es meramente
administrativo: las evaluaciones de latencia de sistemas ROS~2 han mostrado
que los resultados dependen de la configuración de software y hardware de
maneras no obvias \cite{betz2023a, kronauer2021}. El código fuente y la
configuración del benchmark están bajo control de versiones. Cualquier
desviación de la metodología planificada que se observe durante los
experimentos finales se registra y se informa.
